\documentclass[12pt]{article}
\usepackage[utf8]{inputenc}
\usepackage[T1]{fontenc}
\usepackage{hyperref}
\title{Xanagu\'{\i}a Zapotec Language Resources for K-12 Curriculum Development}
\author{Compiled by Claude Sonnet 5.5 (claude-sonnet-5-5) for Andr\'as Kornai}
\date{}
\begin{document}
\maketitle

\section{Introduction}
Xanagu\'{\i}a Zapotec (ISO 639-3 ztg; Glottocode xana1235; endonym Diidz Z\"e) is a Sierra Sur Zapotec language of Oaxaca, Mexico. Sources give about 2,500 speakers (1990 census). Joshua Project gives 4,000 for the people group. Wikitongues rates it EGIDS 6a (vigorous), and it uses a Latin script. No school materials were found, and the language is not known to be taught in schools.

\section{Primary Educational Resources}
None found in search results.

\section{Linguistic and Reference Materials}
\begin{itemize}
\item Wikipedia article on Xanagu\'{\i}a Zapotec. \url{https://en.wikipedia.org/wiki/Xanagu%C3%ADa_Zapotec}
\item Wikitongues page for ztg. \url{https://wikitongues.org/languages/ztg/}
\item ASJP wordlist entry for Xanaguia Zapotec. \url{https://asjp.clld.org/languages/XANAGUIA_ZAPOTEC}
\end{itemize}

\section{Religious and Community Resources}
\begin{itemize}
\item Joshua Project language page for ztg. \url{https://joshuaproject.net/index.php/languages/ztg}
\item Joshua Project people-group page for the Xanaguia Zapoteco. \url{https://joshuaproject.net/people_groups/14606/MX}
\item MegaVoice audio Bible: Gospel of Luke in Zapoteco de Xanagu\'{\i}a. \url{https://megavoice.com/media-cloud/m027249-gospel-of-luke-ztg-zapoteco-de-xanaguia-audio-bible/}
\item MegaVoice audio: Good News in Zapoteco de Xanagu\'{\i}a. \url{https://megavoice.com/media-cloud/m027248-good-news-ztg-zapoteco-de-xanaguia-audio-bible/}
\end{itemize}

\section{Digital and Audio Resources}
\begin{itemize}
\item Global Recordings Network: audio Bible stories in Xanaguia Zapotec (40 stories from Genesis through Acts). \url{https://www.globalrecordings.net/en/language/ztg}
\end{itemize}

\section{Curriculum Development Context}
The language is still actively spoken, but the searches found no primers, dictionaries or school programmes. The available audio is religious in content. A curriculum would need new materials developed with the community, building on the Latin-script usage noted by Wikitongues.

\end{document}